\subsection{Metatheory}
\label{sec:metatheory}

Value well-formedness (\figref{well-formed-values}) determines the types at which an ordinary value is
well-formed, and the context entry of a module or class value; environment well-formedness applies it pointwise,
relating an environment to a context. The predefined environments are well-formed by construction, and a module loads to an environment
well-formed at its signature.

\begin{theorem}[Loaded environments are well-formed]
\label{thm:loaded}
If $\vdash_{\mathsf{load}} q : \Gamma$ then every type in a context entry of $\Gamma$ is well-formed in the
context in which $q$ is checked.
If $\mathcal{P}(q) = (\rho, \Gamma)$ then $\vdash \rho : \Gamma$.
If $\vdash_{\mathsf{load}} q : \Gamma$ and $q \Rightarrow_{\mathsf{load}} \val{\rho}$ then $\vdash \rho : \Gamma$.
\end{theorem}

\begin{lemma}[Variables are bound at run time]
\label{lem:variables-bound}
If $\vdash \rho : \Gamma$ and $\Sigma; \Gamma \vdash x \synth \tau$ then $\rho(x) \neq \Unbound$.
\end{lemma}

\noindent The \ruleName{var} rule requires $x$ to be definitely assigned, and no value other than $\Unbound$ is
well-formed at $\Unbound$ or $\DU{\tau}$, so a well-formed program never raises Python's
\pyinline{UnboundLocalError}. \todo{Unproven.}

